\chapter{Torque}

\section*{Introduction}

Every time you turn a doorknob, pedal a bicycle, or use a screwdriver, you are applying torque. The human body is a torque-producing machine: your muscles generate torques at joints to produce movement. Engineers design engines by maximizing torque output; athletes maximize torque to throw farther and hit harder.

\section*{Fundamental Equation Used}

\[
\boldsymbol{\tau} = \mathbf{r} \times \mathbf{F}
\]
\section*{Assumptions}

\textbf{Assumptions:} The force $\mathbf{F}$ acts at position $\mathbf{r}$ from the pivot point. The force is applied instantaneously. The system is rigid (no deformation).


\textbf{Limitations:} For deformable bodies, internal torques and stress distributions must be considered. At relativistic speeds, torque must be reformulated using 4-vectors. Does not include magnetic torques on current loops (add $\boldsymbol{\mu} \times \mathbf{B}$).


\section*{Mathematical Derivation}

\textbf{Step 1: The Intuitive Origin of Torque.}

Consider a rigid body that can rotate about a fixed pivot point $O$. A force $\mathbf{F}$ is applied at a point $P$ whose position relative to $O$ is $\mathbf{r}$. We want to determine how effective this force is at producing rotation. Intuitively, only the component of $\mathbf{F}$ perpendicular to $\mathbf{r}$ contributes to rotation---the parallel component simply pushes or pulls on the pivot.

\textbf{Step 2: Decomposing the Force.}

Decompose $\mathbf{F}$ into components parallel and perpendicular to $\mathbf{r}$:

\begin{equation}
\mathbf{F} = \mathbf{F}_\parallel + \mathbf{F}_\perp, \qquad F_\parallel = F\cos\theta, \qquad F_\perp = F\sin\theta,
\end{equation}

where $\theta$ is the angle between $\mathbf{r}$ and $\mathbf{F}$. The parallel component produces no rotation (it acts along the lever arm). The perpendicular component $F_\perp = F\sin\theta$ is the only part that causes angular acceleration.

\textbf{Step 3: Defining the Magnitude of Torque.}

The effectiveness of the perpendicular force depends on the distance $r$ from the pivot. The moment arm is $r\sin\theta$ (the perpendicular distance from the pivot to the line of action of $\mathbf{F}$). The magnitude of the torque is:

\begin{equation}
\tau = r\,F\sin\theta.
\end{equation}

This is the product of the lever arm ($r$) and the tangential component of force ($F\sin\theta$).

\textbf{Step 4: The Direction of Torque.}

Rotation occurs \emph{about} an axis, not along a direction. The axis of rotation is perpendicular to the plane containing $\mathbf{r}$ and $\mathbf{F}$. The right-hand rule determines the direction: curl the fingers of your right hand from $\mathbf{r}$ toward $\mathbf{F}$; your thumb points in the direction of the torque vector. This direction convention is exactly captured by the \textbf{cross product}.

\textbf{Step 5: The Cross Product Definition.}

The torque vector is defined as:

\begin{equation}
\boxed{\boldsymbol{\tau} = \mathbf{r} \times \mathbf{F}.}
\end{equation}

In component form, using the determinant representation:

\begin{equation}
\boldsymbol{\tau} = \begin{vmatrix} \mathbf{\hat{i}} & \mathbf{\hat{j}} & \mathbf{\hat{k}} \\ r_x & r_y & r_z \\ F_x & F_y & F_z \end{vmatrix} = (r_y F_z - r_z F_y)\,\mathbf{\hat{i}} - (r_x F_z - r_z F_x)\,\mathbf{\hat{j}} + (r_x F_y - r_y F_x)\,\mathbf{\hat{k}}.
\end{equation}

The magnitude is $|\boldsymbol{\tau}| = |\mathbf{r}||\mathbf{F}|\sin\theta$, consistent with our intuitive derivation.

\textbf{Step 6: Torque as the Rate of Change of Angular Momentum.}

The angular momentum of a particle is $\mathbf{L} = \mathbf{r} \times \mathbf{p} = \mathbf{r} \times (m\mathbf{v})$. Differentiating with respect to time:

\begin{equation}
\frac{d\mathbf{L}}{dt} = \frac{d\mathbf{r}}{dt} \times \mathbf{p} + \mathbf{r} \times \frac{d\mathbf{p}}{dt} = \mathbf{v} \times (m\mathbf{v}) + \mathbf{r} \times \mathbf{F}.
\end{equation}

The first term vanishes because $\mathbf{v} \times \mathbf{v} = 0$. Therefore:

\begin{equation}
\boxed{\boldsymbol{\tau} = \frac{d\mathbf{L}}{dt}.}
\end{equation}

This is the rotational analogue of Newton's second law ($\mathbf{F} = d\mathbf{p}/dt$). For a rigid body with moment of inertia $I$ rotating about a fixed axis, $\mathbf{L} = I\boldsymbol{\omega}$, so $\boldsymbol{\tau} = I\boldsymbol{\alpha}$, where $\boldsymbol{\alpha} = d\boldsymbol{\omega}/dt$ is the angular acceleration.

\textbf{Step 7: Static Equilibrium.}

A rigid body is in rotational static equilibrium when the net torque vanishes:

\begin{equation}
\sum_i \boldsymbol{\tau}_i = \sum_i \mathbf{r}_i \times \mathbf{F}_i = \mathbf{0}.
\end{equation}

Together with the translational equilibrium condition $\sum \mathbf{F}_i = \mathbf{0}$, these equations fully determine the statics of rigid bodies.

\section*{Characteristics of the Equation}

Torque is a pseudovector (changes sign under mirror reflection). The net torque equals the rate of change of angular momentum: $\boldsymbol{\tau} = d\mathbf{L}/dt$. In static equilibrium: $\sum \boldsymbol{\tau} = 0$. Torque has SI units of N$\cdot$m (newton-meters).
\section*{Solution Methods}

Free body diagrams with torque analysis, choice of pivot point to eliminate unknown forces, system of simultaneous equations for static equilibrium, $\tau = I\alpha$ for rotational dynamics.
\section*{Verifications}

Torque sensors measure applied torques directly. Balance experiments: when $\sum \tau = 0$, a system is in rotational equilibrium. High-speed cameras can measure angular acceleration and infer torque.
\section*{Applications}

Wrench and lever design, gear systems, bicycle mechanics, automotive engines, helicopter rotors, door hinges, crane stability, human biomechanics (joint torques in walking and running).
\section*{What Richard Feynman Would Have Asked}

\subsection*{1. The Plain-English Translation}
\textit{``What is torque really? Not the formula---the concept.'''}

The Core Meaning: Torque is ``twisting force.'' Just as force causes linear acceleration, torque causes angular acceleration. But torque also depends on where you apply the force---the same force applied near the pivot produces less torque than the same force applied far from the pivot. This ``where'' is just as important as ``how much.''

The Metaphor: Imagine pushing a merry-go-round. Pushing near the center barely rotates it (low torque). Pushing at the rim rotates it easily (high torque). Pushing tangent to the rim at the farthest point gives maximum torque. The formula $\tau = rF\sin\theta$ captures all of this: $r$ is where, $F$ is how hard, and $\theta$ is the direction.

\subsection*{2. The Localized Mechanics}
\textit{``At the point where the force is applied, what is happening to the material?''}

He would press you on the stress distribution. When you apply torque to a bolt, the wrench exerts a force on the bolt's surface. This creates a stress field inside the bolt---shear stresses near the surface that decrease toward the center. The bolt rotates when the torque exceeds the friction torque holding it in place.

The Smoothness Check: He would ask why torque is a cross product rather than a dot product. Because torque has direction---it produces rotation \emph{about} an axis, not along a direction. The cross product naturally gives a vector perpendicular to both $\mathbf{r}$ and $\mathbf{F}$, which is the axis of rotation.

\subsection*{3. Testing the Extreme Limits}
\textit{``What happens at extreme torques?''}

The Zero-Torque Limit: When $\tau = 0$, angular momentum is conserved. This is why a figure skater spins faster when pulling in their arms---smaller moment of inertia, same angular momentum, higher angular velocity.

The Infinite-Torque Limit: In practice, infinite torque doesn't exist. But very large torques can deform or break materials. The yield torque is the maximum torque a shaft can transmit before permanent deformation occurs.

The Gravitational Torque: Tidal forces on the Moon create a torque that has synchronized the Moon's rotation with its orbital period (tidal locking). The same mechanism is gradually slowing Earth's rotation and will eventually lock the Earth-Moon system.

\subsection*{4. Dimensionality and Geometry}
\textit{``Does the shape of the object affect how torque works?''}

He would point out that the moment of inertia $I$---which relates torque to angular acceleration ($\tau = I\alpha$)---depends entirely on the shape and mass distribution. A hollow cylinder has more $I$ than a solid cylinder of the same mass (mass is farther from the axis). This is why flywheels are designed to concentrate mass at the rim.

\subsection*{5. Cross-Disciplinary Connection}
\textit{``Where else does the concept of torque appear?''}

Torque appears in: magnetic dipole moments ($\boldsymbol{\tau} = \boldsymbol{\mu} \times \mathbf{B}$---how compass needles align), molecular chemistry (bond angles and rotational barriers), economics (``leverage''---small forces at strategic points producing large effects), and politics (``political leverage''---applying influence at the right point to produce disproportionate change).
\section*{FAQ}

\textbf{Q: Why does a longer wrench make it easier to loosen a bolt?}\\
\textbf{A:} Because $\tau = rF\sin\theta$. Increasing $r$ (the lever arm) increases the torque for the same force. A wrench twice as long produces twice the torque, requiring half the force.

\textbf{Q: Can torque be negative?}\\
\textbf{A:} The sign indicates direction (clockwise vs counterclockwise). Conventionally, counterclockwise torque is positive. A negative torque produces clockwise rotation.
\section*{Important Bibliography}

\begin{enumerate}
\item Morin, D., \textit{Introduction to Classical Mechanics}, 2nd ed. (Cambridge University Press, 2008), Chapters 6--7.
\item Marion, J.B. and Thornton, S.T., \textit{Classical Dynamics of Particles and Systems}, 5th ed. (Brooks/Cole, 2004), Chapters 10--11.
\item Symon, K.R., \textit{Mechanics}, 3rd ed. (Addison-Wesley, 1971), Chapters 7--9.
\item Meriam, J.L. and Kraige, L.G., \textit{Engineering Mechanics: Statics}, 8th ed. (Wiley, 2016), Chapter 3.
\end{enumerate}


